\section{Beau Geste at Notre Dame}

On 21 May 2013, the new right, formerly old right, French historian \textbf{Dominque Venner} shot himself in the mouth using a Belgian 9mm pistol in front of the altar at Notre Dame Cathedral in Paris. M. Venner considered his final act as a dual protest against the new French law allowing ``gay marriage'' and the ``replacement'' of the French people by immigrants from Africa and the Maghreb. \textbf{Friedrich Nietzsche} passed the breaking point into insanity when he witnessed a horse being beaten. Dominique Venner passed that point when he read this in a blog entry by an Algerian, which ``chilled his spine'':

\begin{quotex}
In any event, in 15 years Muslins will be in power in France and will repeal that law.
\end{quotex}


The logic of that is unassailable and led M. Venner into a state of inner contradiction: Either end immigration and retain gay marriage or else increase immigration and gay marriage will disappear. This existential angst led him to his ostentatious ``beau geste''. This is how he justified himself in a blog entry from earlier that day:

\begin{quotationx}
Certainly new, spectacular, and symbolic gestures are necessary to shake up the lethargic, to rattle the consciousness of those anaesthetized, and to awaken the memory of our origins. We are entering a time where words must be authenticated by acts.

It is necessary to also remember, as Heidegger brilliantly formulated it, that the essence of man is in his existence and not in ``another world''. It is here and now that our destiny is played out right to the last moment. And that last moment has as much importance as the rest of one's life. That is why it is necessary to be oneself up until the last instant. It is in deciding oneself, in truly willing one's destiny, that a man is the victor over nothingness. And there is no way out of this requirement since we have only that life in which it is possible either to be completely ourselves or to be nothing.
\end{quotationx}


Now a beau geste appears noble in form, but is meaningless in substance. Nevertheless, there are those attracted to that gest, like vultures hovering over a carcass; they are quick to discern some meaning in that act of despair. Thus as a jest is not funny when you have to explain the punch line, a geste is not meaningful if you have to explain the meaning. Certainly those who were intended to be shocked, shaken, and awakened, derived nothing from it.

\paragraph{Population Replacement}


What M. Venner failed to notice is that population replacement has been going on for quite some time. The France that was shaped over the course of 1000 years by 40 Catholic kings---as Charles Maurras put it---had ended 225 years ago. Its Catholic population, structured by the three estates, was replaced by Jacobins, atheists, philosophes, Masons, Huguenots, and other assorted aliens. Just as a man is ``brain dead'' before the complete lifelessness of the body, so also was France spiritually dead a long time ago; in social terms, it may take a few generations for that death to be fully noticed on the material plane.

Into that existential situation, M. Venner was thrown. His failures reflect his failure to understand it. He flailed about, looking for some spiritual value to revive the corpse. Dissatisfied with what he saw available, he went back 3000, then 30,000 years. He believed life moves in a circle, but the circle did not close; in that respect, it was more like a line and all he could see was the end of the line, not the recurrence of the past.

\paragraph{The God}


Unfortunately, M. Venner was part of the vulgarity he opposed. As we know, the revolutions from below, like those of the Jacobins and the Bolsheviks, are always atheistic; they deny transcendence, and want only to overturn the established order of things. Likewise, for M. Venner all that counts is ``existence''. In that case, his view of existence is no more valid than that of the bus driver who transported him to Notre Dame. So he wanted a beau geste to awaken the masses, but it is the ideology of the masses that is itself the problem. If there is nothing but the existence of individuals, then sexual orientation or country of origin---the primary concerns of M. Venner---has no significance. There is nothing to awaken.

What he could not come to grips with, as he had no Negative Capability\footnote{\url{http://www.gornahoor.net/?p=6305}}, was that the Islam he hated had become the bearer of Tradition in France. He weakly objected that the West, unlike Islam, respects women. If he means contemporary, sexualized and militarized Western women, ready for front-line combat, his case may be hard to make. This is the main point made by \textbf{Guido De Giorgio}\footnote{\url{http://www.gornahoor.net/?p=6420}}.

Ultimately, M. Venner has nothing to offer us men of Tradition. Rather than drawing on France's 1000 year spiritual heritage, he relies on the ``brilliant formulation'' of \textbf{Martin Heidegger}, the German. He pays lip service to some mythical ethnic spiritual tradition that would unite the French, much like the Jews, but forgets that a Tradition is not man-made but descends from above. Rather he should have pulled from France's spiritual tradition what he could, just as it had preserved and integrated the important features of the earlier mystery cults and neo-Platonic systems\footnote{\url{http://www.gornahoor.net/?p=6455}}.

Instead he offers us the vulgarity of the masses dressed up as something higher. The mystery religions taught the One God, just as did the Academy, the Lyceum, and the Stoa. That was unknown to the masses who were polytheists or atheists. Only with the Christian teachings was the God made known to all. The New Right promotes everything opposed to Tradition: atheism, neo-paganism, zoological racism, ethnic anti-Semitism, and even, ironically (considering M. Venner's position), homosexuality. Everything M. Venner found valuable in the glory of Greece and the grandeur of Rome, then so do we. As we've said many times, we love our pagan past. Yet we love our recent past even more.

He considered the ultimate choice between one's identity and death as a brilliant insight. Since he could no longer identify as a Frenchman as it is now understood, he chose death. However, we see things differently. We have to \emph{establish} our identity, to realize (make real) all our possibilities; that is a task in time, not a static state of being. Death will come soon enough. In the meantime, we engage in spiritual combat, both internally and externally. We will win, we will lose, but ultimately we will prevail since the world is a line not a closed circle.

\paragraph{Beau Geste Addendum}

Since today is a family day, I can only offer a few quick thoughts, probably more than any ``fly by'' comments deserve.

As for the current situation in Europe (and the West by extension), there are no surprises for those who understand Tradition. As a matter of fact, one of the very first posts here had to do with some ideas from Rene Guenon on this topic: Please refer to The Prophecies of Rene Guenon\footnote{\url{https://gornahoor.net/?p=35}}. I urge you to consider the options presented carefully; it may make the more recent post clearer.

Comments without any actual content are a waste of my time; nevertheless, since they are now on public display, I will try to provide the missing content.

This is an example of a commentary that can be described as ``petty'': Dominique Venner's Final Solution\footnote{\url{http://www.newyorker.com/online/blogs/newsdesk/2013/05/dominique-venners-final-solution.html}}.

This is one that is disrespectful since it judges the geste by its own standards and allows unsupported calumnies: La provocation de Dominique Venner\footnote{\url{http://yvesdaoudal.hautetfort.com/archive/2013/05/21/la-provocation-de-dominique-venner.html}} (English excerpts here: Thoughts on Dominique Venner\footnote{\url{http://galliawatch.blogspot.com/2013/05/thoughts-on-dominique-venner.html}}).

No one does any service to a true intellectual by according him blind homage. M. Venner gave us a justification for his geste as well as what he hoped to accomplish. To pay him proper respect is to take that seriously. That is what I attempted to do.

M. Venner wanted to ``shake, rattle and awaken''. This is his standard, not mine. Hence, to show respect, I judge him by that standard. I would prefer some comments on that standard rather than snide remarks.

In the light of Guenon's predictions, even had he succeeded, it would be inadequate. Guenon has also written that the end of an era is the end of an illusion. That is a daunting task, but the one thing necessary, and few people want to be reminded of that.

M. Venner came to the only conclusion possible for him: to choose his identity or to choose nothingness. He chose the latter. This is not my opinion, that is what he wrote and what he did. With that, he offers up the same choice to everyone else: choose nothingness or establish your identity. The problem remains is that he never made clear what is our identity, except in a negative sense: our identify is not-Muslim.

M. Guenon offers a different set of choices.


\flrightit{Posted on 2013-05-26 by Cologero}

\begin{center}* * *\end{center}

\begin{footnotesize}
\begin{sffamily}
\texttt{John Morgan on 2013-05-27 at 00:53 said:}

How one chooses to judge and interpret Venner's final act is a matter of opinion (although what has happened to ``De mortuis nihil nisi bonum''?).Nevertheless, this makes a number of factual errors:

1. ``M. Venner considered his final act as a dual protest against ``gay marriage'' and the ``replacement'' of the French people by immigrants from Africa and the Maghreb.'' Venner specifically denied any connection between gay marriage and his suicide in his final Website post: ``The May 26 protestors cannot ignore this reality. Their struggle cannot be limited to the rejection of gay marriage. The ``great replacement'' of the population of France and Europe, denounced by the writer Renaud Camus, is a far more catastrophic danger for the future.''

2. ``He flailed about, looking for some spiritual value to revive the corpse. Dissatisfied with what he saw available, he went back 3000, then 30,000 years.'' This makes Venner out to be some sort of neo-pagan, which isn't the case. He never spoke ill of the Catholic Church or advocated for paganism to replace it. His criticism of it was that, in modern times, the Church embraced Cartesian rationalism and liberalism, a criticism he shared with the traditionalists. Also, Venner never attributed the decline of the West entirely to a lack of ``spiritual values'' -- his concept of ``traditionism'' was not a metaphysical or religious one as with Evola and Guenon, but was how he designated the unique history, identity and culture of a people, as well as its religion. So to say that Venner was looking merely to revive ``spiritual values'' is not correct, since for him that would be insufficient on its own.

3. ``Rather than drawing on France's 1000 year spiritual heritage, he relies on the ``brilliant formulation'' of Martin Heidegger, the German.'' Venner always denied being merely a French nationalist, and maintained that he saw himself as a European who happened to be French. By ancestry, he was actually Celtic and German, not French. So there's no inherent contradiction in that, as this statement seems to be implying.

4. ``The New Right promotes everything opposed to Tradition: atheism, neo-paganism, zoological racism, ethnic anti-Semitism, and even, ironically (considering M. Venner's position), homosexuality.'' This overlooks the fact that Venner was only involved with the ``New Right'' during its first decade of existence. By the 1980s he was no longer involved with Benoist's GRECE. So Venner was not, strictly speaking, a ``New Right'' author. Beyond that issue, however, while I would agree that the New Right is not a traditionalist movement in the sense of Evola or Guenon, one cannot accurately claim that it has ever stood for ``atheism'' (where have Benoist or any of the others advocated for that? Quite the opposite), ``zoological racism'' (the New Right has always stood against racism), or ``homosexuality'' (again, I have never come across a single place where the New Right has defended it). Even ``neo-paganism'' is debatable, since the New Right has never argued that Westerners should take up the worship of Odin or Jupiter again, as Benoist makes clear in his book on paganism, where he calls for a more ``philosophical paganism.''

\hfill

\texttt{Constantine on 2013-05-27 at 01:03 said:}

It is about time for a conference, congress, meeting or what ever you want to call it. It must take place soon in order, as you say establish an identity, plus the issue of governance and finally how to respond against these forces of darkness.

(Matthew 13:31-32)

\hfill

\texttt{Ash on 2013-05-27 at 02:37 said:}

I'm afraid I must disagree on your comments about Islam. By and large, the Islam being promoted by the fundamentalists in Europe, the street patrols of Britain and so on, is of a degenerate Wahhabist/Salafist type at its most theoretical. It has as much tie to the sacred Tradition as the megachurch Evangelicals. In practice, it becomes a political ideology, resorting to mob rule and having no knowledge of Shariah beyond the call to implement it.

Having read his final post and suicide letter, in addition to other writings of his, I don't believe that he thought the choice was between identity and suicide.

``It is here and now that our destiny is played out right to the last moment. And that last moment has as much importance as the rest of one's life. That is why it is necessary to be oneself up until the last instant. It is in deciding oneself, in truly willing one's destiny, that a man is the victor over nothingness.''

What I gather from his quotation of Heidegger, Venner believed that his death was an act which gave expression to his identity as a Frenchman and European in communion with the European tradition (as he understood it). It was not a suicide out of fear or abandonment. His note gives a clear impression that his death is a choice made in sound mind and will. Venner's last book is on the samurai and their tradition, something anyone on Gornahoor should be familiar with. It seems likely that this influenced his decision greatly. Doubtless, his understanding of Europe's tradition was flawed. But then, I doubt any man but the Saints who have seen God can claim a pure understanding. I can speak only for myself, but I will not judge the inner and spiritual state of such a man, but rather hear his call and heed it by coming back to the tradition France and Europe has abandoned, even as we walk in the current age towards uncertain times ahead.

\hfill

\texttt{Avery Morrow on 2013-05-27 at 02:57 said:}

90\% of the responses to metaphysics fall into the categories of ignorance, misunderstanding, or misapplication. M. Venner, if this characterization is correct, was in a state of ignorance, seeing good in his ancestry but not knowing the relationship between that and present society. Misunderstanding of what one has read is quite visible on Internet blogs, and even a precise and accurate knowledge is so easy to misapply that Schuon did it.

I agree that the time is ripe for an in-person congress --- the Internet seems to foster misunderstanding in particular, because it's so easy to sound reasonable when you are only talking to yourself.

\hfill

\texttt{Mihai on 2013-05-27 at 03:45 said:}

I agree with almost everything you said here.

Certainly, the new-right is absolutely useless and completely devoid of genuine principles. I do believe that should an ideology such as that of the European and American new-right come to power it will only aid the advancement and even the fulfillment of subversion.

As far as I can tell, new-rightists seem to believe that they are `traditional' and oppose modernity because they support hierarchy and oppose democracy.

Such things are completely irrelevant. Hell also has its hierarchy (or downarchy- as CS Lewis put it).

\hfill

\texttt{John Morgan on 2013-05-27 at 04:33 said:}

The fact that you mention the New Right ``coming to power'' shows that you don't really understand what it is. The European New Right has always been a metapolitical school of thought, not a political movement. There has never been a New Right political party that fights in elections or advocates for revolution. As for New Rightists believing they are traditionalists, while Alain de Benoist and others in GRECE have occasionally made reference to Guenon, Evola, and so forth, Benoist has always freely admitted the fact that his is not a traditionalist school (which would abrogate the need for metapolitics altogether). It seems to me that you are confusing the European New Right with the movements that go by that name in the US and the UK, which share no organizational and little philosophical or political ideology with the European movement. Benoist, in fact, has explicitly rejected any association between GRECE and the Anglo-American ``New Right'' movements.

All of this is besides the fact that Venner himself was not really a ``New Right'' author, as I detailed in my own response to this article.

\hfill

\texttt{V.O. on 2013-05-27 at 05:17 said:}

What is the potentially metaphysical significance of ethnos or racism? It would be a means by which to preserve the homogeneity of the given tradition, supporting its transcendent orietnation despite being in the world; preventing syncretism, corruption of its principles etc. According to the wordliness of man, a sense of ethnos becomes almost necessary to support the purity of the ``metaphysics.''

In societies with a strong ethnos, it is obvious that the tradition remains relevant, even if at times dormant. A Russian, Serbian or Greek that would turn away from Orthodoxy, Pole or Croatian turning away from Catholicism or Palestinian turning away from Islam has an almost pathological connotation. In this sense dissolution of the ethnos is also an antichrist tactic for dissolution of the tradition, encouraging syncretism, spiritual tourism, etc. Modern world's elision of ethnos corresponds to the elision of tradition.

Satanism of course is when the vertical dimension for which the ethnos exists is omitted in favour of only the worldliness of the ethnos. Again the Latin proverb that Upton among others cites frequently shows its depth: corruptio optimi pessima. Ethnos as a form of life within the world that, as it were, defends the integrity of the given tradition degenerates into a vulgar and populist materialist racism.

\hfill

\texttt{Mihai on 2013-05-27 at 05:28 said:}

The problem, V.O., is that ethnos descends, organically, from tradition, not the other way around. The lack of tradition and atheism are not results of globalism, but globalism is a result of the two.

The new-right surely does not understand this, and so relegates ethnos, like you said, to a purely bi-dimensional, materialist affair.

\hfill

\texttt{V.O. on 2013-05-27 at 05:53 said:}

Whereas ethnos follows from tradition, there is still a ``relative autonomy'' in the world -- if we wanted to use Marxist-Althusserian terms it is a ``determination in the last instance'', but better to use terms such as ``natural will'' and ``gnomic will.'' The natural will is this relative autonomy which is a true freedom, whereas the gnomic will as recently discussed upon on this site, is reduced to the ``free choice''. Freedom and choice are not the same, Freedom retroactively negates the ``choice'' and transforms the cosmos -- i.e., the first Adam's decision for death. to paraphrase St. Cyril of Alexandria: Adam's decision for death is a universal death.

Free choice-the gnomic will aims to preserve the form of the ``choice'' in the future ``consumerism'', keeping the world as it is; freedom in contrast always transmogrifies the world.

The reason for my preamble here Mihai is in response to globalism as a result of this lack of tradition -- the foundation of your account here is the traditional-Neo-Platonistic account of privation, whereby globalism or the modern world is simply the consequence of this privation. But the principle of privation is not what is at stake in Christianity -- there is the relative autonomy or freedom of the world, and thus anti-traditional forces are certainly active in the world -- their tactics are the exercising of this same freedom, as Adam chooses death.

The ethnos is therefore a reflection of this freedom as opposed to choice -- it is an expression of the natural will as opposed to the gnomic will. Although I live in a Slavic non-Orthodox country, I am still Russian and therefore there is no option for me other than Orthodoxy -- this is the freedom that is not a choice, There simply is no choice, and one of the functions of ethnos -- in addition to its decisive metaphysical function as contributing to the preservation of the integrity of the given tradition -- is to help maintain the separation of freedom from its modern-capitalsit--consumerist vulgarization as choice/gnomic will.

Ethnos, of course, is not necessary to these acts, just as the creation itself is unnecessary -- but from the metaphysical perspective towards the world, from the perspective of God's oikonomos to the world, it does have a positive effect, although one that, as we have noted, can quickly turn into its opposite when the transcendent is omitted.

\hfill

\texttt{Izak on 2013-05-27 at 13:04 said:}

I appreciate the vast majority of what Gornahoor has done since it started, but I agree with your response, and I'm glad that you wrote a much better one that I could hope to write.

I found this post here to be disrespectful and petty. I have nothing else to add.

\hfill

\texttt{Jason-Adam on 2013-05-27 at 16:54 said:}

While I oppose all suicide as cowardice and immoral, I can not but find myself in agreement with M. Venner.

Islam may be a tradition but is not MY Tradition. As de Giorgio eloquently wrote, ``paganism'' and Catholicism are two eras of the singular Roman Tradition. To become a Muslim -- to become ``oriental'' as Evola put it -- is to abandon the heritage of my ancestors. I can not live that way. I speak as someone who has prayed in mosques and took classes with the intent of converting to Islam but it didn't work for me -- it's just too alien.

I actually am a ``globalist'' though not of the capitalist bourgeois American kind, I believe in one Roman Empire covering the entire globe -- order will come when one ruler replaces many.

Part of rejecting egalitarianism means rejecting the notion that all religions and races are equal -- I am proud to say I believe that European man (speaking of a personality type not a physical body) is the superior race destined to rule the Earth if only he had kept to his Tradition. I say this not out of pseudoscientic racism of a Chamberlain nor a naturalism of a Nietzsche. I say it because my Church which is Roman is the One True Church to which all men must belong to for salvation from hell.

Julius Evola offers a way to resolve the psychological turmoil that destroyed M. Venner -- we must understand that the world is ending, there is no hope for us at all politically and we must be strong and live as isolated from the world as we can. Our citizenship is in the ideal republic, not the degenerate Athens of our day.

Desire for a just political order is another craving we must abandon to achieve deliverance. The true initiate does not desire to rule, he is forced to be a ruler as Plato taught us.

\hfill

\texttt{Ash on 2013-05-27 at 17:20 said:}

Then surely the task at hand is to take what is true or valid in the New Right and contextualize it within the teachings of Tradition? ``Everything that is good, is ours.'' It is always better to educate people and win them to your side (without sacrificing truth or honour, of course) than to force a situation of opposition.

\hfill

\texttt{Graham on 2013-05-27 at 17:54 said:}

That's well-stated, V.O.

\hfill

\texttt{Cavalcare on 2013-05-27 at 18:27 said:}

Izak, you are being very generous my friend. Maybe not this blog (although this is arguable, considering the direction it's been taking), but certainly Cologero's legitimacy has officially descended to below ground level with this post.

\hfill

\texttt{Pickman on 2013-05-27 at 23:58 said:}

Of course, modern Islam in the west today is entirely degenerated into a cultural protest movement of mob rule that is largely supported by hard-left and socialist opportunism. As stated above it has deteriorated to such an extent that Shariah Law is a necessary act of petty enforcement upon its supporters. Even in the height of the Christian middle ages the peoples of the cities were trusted with basic freedoms and women within our societies were not treated too high or too low (like in Islam).

Little to none of its spiritual ``grace'' among its supporters can be barely perceived and Sufism is outlawed amongst mainstream Islam (Sunni or Shia, Wahhabist/Salafist), therefore the Darkness will only grow if it is allowed to fester. The author would do well to dispel of illusions and watch it happening from the ground up rather than treating Islams ``integration'' into European societies as something to be welcomed. It should be noted that neither Evola or Guenon witnessed these current events and had the envious privilege of studying said Islamic societies from the plexiglass like they were just another harmless exotic Asiatic tradition. This is despite Guenons later conversion -- just a pity he was denied the Indian passport... Some will say Venner was a Martyr others have likened his death to Roman Devotio -- which indeed has more noble sentiments. Whatever the case, he was a brave and loyal man, and will not have to suffer the indignant death of rotting in some decrepit state run nursing home swamped by Asian/African illegals.

\hfill

\texttt{Pickman on 2013-05-28 at 00:06 said:}

Exactly, well said. The writings of New Rightism on Gornahoor are often greatly uninformed and often sounds like it was lifted from a searchlight/Splc report.

\hfill

\texttt{Pickman on 2013-05-28 at 00:16 said:}

Entirely true. Modern Islam in the West is a degenerate hybrid of deracinated uncouth mobs who exist as a fundamentalist protest movement -- many of which share stark political bedfellows with socialist and Hard-left currents. Something Guenon would even have disproved of. Sufism is outlawed amongst the most mainstream of these sects and no light will ever come from such darkness, much less, the ridiculous notion that it can save the West from itself (no matter how many centuries pass).

\hfill

\texttt{Mihai on 2013-05-28 at 03:36 said:}

I didn't mention the new-right coming to power, I mentioned a group upholding such an ideology as that upheld by the new-right coming to power, which I don't regard as so improbable given the latest developments in the western countries.

Metapolitical or not, the new-right is intellectually bankrupt, its ideology amounts to little more than what in the religious field we would call syncretism. Its ideology is sure to be taken as the support for the future phase of subversion.

\hfill

\texttt{Mihai on 2013-05-28 at 04:10 said:}

John Morgan,

I appreciate Arktos for some of the high-quality publishing of some important books. The fact is, however, that the books I appreciate and which can be found on that website, completely contradict the very postulates of the new-right.

The books that are specifically new-rightist differ in very little from the ideological salad promoted by most post-modernist currents. Benoist's book on paganism is a very clear example. If Benoist believes that the reconstruction

of Europe is to find a base in ``the return'' to the pagan roots (philosophical or otherwise), is enough to prove that he and his movement have received more attention on this blog than they deserve. If Nietzsche is his primary source of inspiration we can lay back and enjoy the continuing of the fall.

\hfill

\texttt{John Morgan on 2013-05-28 at 05:36 said:}

Dear Mihai,

I was not attempting to defend the New Right on traditionalist grounds. As I wrote in my other response to you, it is clearly not a traditional movement, nor has it ever claimed to be such. Traditionalists can take it or leave it as they please. However, that does not excuse the gross mischaracterization of the movement given in the original essay, nor its other factual errors. That was the only purpose of my reply.

Thanks for supporting Arktos, however.

--John

\hfill

\texttt{Pickman on 2013-05-28 at 07:36 said:}

Indeed, there is one thing to constructively criticise a Rightist cause which is still in its gestation period and not fully matured (such as what has been labelled the ``New Right'') but it is quite another to spite it and mock it through parody which is unfortunately what Gornahoor does. This (left-like deconstructive) tone is unfair I believe and frankly unwarranted -- especially to potential sympathisers. I'll note in fair dues that Cologero has contributed articles to Counter-Currents and has provided us (the readers of Gornahoor) with some excellent translations and articles on traditional subjects which keeps me returning to this site -- however it's unhelpful to throw stones in glass houses.

It was once said on here that the Rights differences are legion but the Left moves in unison. Another trouble is that the Left never negates or scathingly attacks its own but the Right does all the time in the hope of winning some respect or appeasement points from its enemies. Again another delusion.

By the way, my posts above may seem like duplicates in content due to the fact that I was unaware they were awaiting moderation and thus thought they had been deleted, hence the repeats.

\hfill

\texttt{Cologero on 2013-05-28 at 08:36 said:}

Pickman, Izak, John Morgan and others:

I appreciate the time you have taken to provide this feedback. I want to apologize for any confusion caused. The post was really about two different topics, one on M. Venner's suicide, the other from notes for an upcoming review of some of Evola's books from Arktos. By combining the two ideas, the tone came across poorly and made it seem like a personal criticism of M. Venner. That was not my intent.

\hfill

\texttt{John Morgan on 2013-05-28 at 09:52 said:}

Dear Cologero,

Thanks for responding. It is true that I interpreted what you wrote as a criticism of Monsieur Venner personally. It's difficult to interpret phrases such as '' M. Venner was part of the vulgarity he opposed'' in any other way. But if that was a mistaken impression, thanks for clarifying it. As I said, however, my own problem wasn't primarily the tone of the piece, but simply the mischaracterizations in it. As I already wrote, Venner could only be termed a ``New Right'' author in a very incomplete sense, considering that his intellectual career began more than a decade before GRECE was formed, and that his association with the group ended about 35 years ago, thus comprising only a small part of his overall intellectual career.

On a related note, if this essay is indicative of what you are going to write in your reviews of the Arktos translations of Evola, I feel I should clarify that Arktos is not a ``New Right'' publisher. We have no overarching ideological mission or identification, although of course we gravitate towards certain themes and issues (including traditionalism and the New Right) that fall under a very broad rubric. But of the more than 50 titles we have published to date, there are only 7 that were written by authors affiliated with the ``New Right'' in any way.

--John

\hfill

\texttt{Cologero on 2013-05-28 at 10:54 said:}

You're correct, John, that attempt to link the ideas failed, since they are really unrelated. The topic will actually be Evola's statement that the state creates the people, not unrelated in some ways to Benoist's distinction the ethnos and the demos. The books stand on their own and nothing I've written should be taken as an attempt to pigeon hole or criticize Arktos in any way.


\hfill

\texttt{Perennial on 2013-06-11 at 01:52 said:}

It may just be me, but it seem like some of the responses are getting rather whiney. The posts I am reading are all ``I feel'' and ``this is a mischaracterization'' without really saying why. I hope for more substance in the future, and less declarative statement. ``Where's the beef?''

\hfill

\texttt{Perennial on 2013-06-11 at 01:53 said:}

Excuse me, I meant to say ``comments'' not posts, as this was not aimed at Cologero for whom I found no fault in this essay.

\hfill

\texttt{White Rabbit on 2014-02-09 at 00:51 said:}

Why is ``ethnic anti-Semitism'' opposed to Tradition?

\hfill

\texttt{Cologero on 2014-02-09 at 23:47 said:}

White Rabbit, ordinarily I would dismiss your question as mere trolling. However, I was really impressed, and somewhat flattered, that you would read old posts on Gornahoor late on a Saturday night. Your desire for intellectual and spiritual development would bring a lump to my throat were I still capable of such emotivity.

As I am sure you understand, there is always a battle of worldviews, which are often opaque to each other. The differences in fundamental assumptions make conversation difficult without good faith efforts. Hence, mere mantras never have the intended effect. Besides, they can only serve the majority of men whose entire worldview consists of a few sentences, each of which can fit on a placard.

The Traditional worldview is a matter of ``vision''; one learns to ``see'' that it is true. That requires and intellectual conversion, i.e., a fundamental reorientation of how one experiences, views, reacts to, and comprehends the world. Although some blessed few may be converted spontaneously, for most of us it requires and effort. The first effort is to clear out one's mind of all mere opinions and think ``as if'' the traditional worldview is correct. Perhaps, then, something significant may result. However, that requires some faith and a decision, which is not really so easy.

The fundamental ``as if'' is that the ``subtle rules the dense'', i.e., spirit, not matter, is fundamental. Hence, we strive to view the world in that way. As far as races go, then, we are more interested in ``spiritual races'', which determine biological races, not the other way around. There is no precise mapping from spiritual races to what most people would regard as race.

As for Jews in particular, they do not appear to be a ``race''. That is what Evola himself claimed. On the scientific and anthropological front, the recent book by Shlomo Sands tries to prove the same thesis. The Medieval period of Europe, which was a traditional civilization, did not accept the idea of zoological antisemitism. So that is the explanation in the nutshell, if you decide to explore things further.

In a planned series on the foundation of the Middle Ages, I will write about this topic in more detail, as well as the idea of nationalism. At the medieval period, these issues were understood as spiritual categories. Unfortunately, over time that understanding was lost and replaced by a more material and physical understanding. That led to a lot of mischief in the world, even under the guise of tradition. That makes conversation impossible as it immediately arouses suspicions. ``Rash judgment'' is a serious moral flaw, so we must take care to understand individuals and groups based on their dignity as persons, and not as mere biological specimens.

\hfill

\texttt{X on 2014-02-10 at 00:25 said:}

Excellent reply, Mr. Salvo.

\texttt{Jason-Adam on 2013-05-27 at 21:37 said:}

My identity is certain as well -- I am a knight of Christendom. Even if Christendom is dead I remain en guarde and ready for battle. I prefer to live and die alone and unknown rather than surrender... ... ... Most people whom I speak to who aren't read up on Tradition think

of me as an apocalyptic nihilist, I simply respond I await the end of the world... ..it is coming... ..Kali Yuga is ending... ... Christ regnabit... ... ... .

\hfill

\texttt{Michael on 2013-05-27 at 22:55 said:}

Jason-Adam, what kind of a knight are you? Which order?

Cologero, I think Guenon's second option would result in the annihilation of Europe. It might be replaced by a culture that is more traditional in some ways, but it would represent a gigantic loss over the European Middle Ages. What a fall!

That leaves ``transformation'' as the only desirable option.

The question is how to achieve this transformation. I'd be interested to read what can be done besides riding the tiger.

\hfill

\texttt{Pickman on 2013-05-27 at 23:21 said:}

Let it be said that Guenon's abdication to Islam was a symbolic death within itself and is simply unacceptable to those who do not wish to convert to an oriental soul which is not congruent to the special character of the Traditional Roman Hyper-Borean Aryan spirit.

The indigenous European people and their western counterparts may be destroyed through degeneracy and miscegenation but No Semitic Tradition born of uncouth slave masters under dhimmitude will save it despite any idealistic illusions those have of Islam and Muslims. Their southern oriental soul and vulgar spirit of their race; worshipers of the feminine and lunar crescent; eternal enemies of Europe's ancestors is not who we are or what should we become. Our divine destiny will be realized by a return/remembrance to said Aryan spirit last seen in classical Romanity in order to stand a chance or perish.

\hfill

\texttt{Logres on 2013-05-28 at 00:26 said:}

What if his suicide has the opposite effect intended? What if it further demoralizes those it is meant to awaken? I wonder if he considered this.

If Europe goes Islamic, it will be an analogous, and greater loss, than when Luther lead it into Protestantism. See Oswald Spengler's comparisons between Luther and Muhammed.

\hfill

\texttt{Izak on 2013-05-28 at 01:26 said:}

I had to hesitate before I used the word ``petty'' because I respect this blog very deeply and the work you do, Cologero.

But John Morgan's criticisms were absolutely correct, and ultimately I stick by my choice of words. You did not respect Venner because you mischaracterized him on several levels, and you shifted your criticisms toward Venner into a critique of something called ``The New Right'' which apparently is both everything and nothing all at once in everyone's mind.

When a man dies, that alone is no excuse to use the occasion to start attacking the entire umbrella under which he stands and everyone peripherally associated with it. If you do not understand why this might look petty to some of your readers, I simply cannot help you. If you wanted to judge his suicide on its own merits, you had no business acting as though his suicide was meant to represent ``The New Right'' and not simply him as a man. It seems to me, and correct me if I am wrong, that you figured, ``Well, he was influenced by Heidegger, and he had some associations with Alain de Benoist, and the New Right is basically Heidegger-influenced, so I will explain everything that I find wrong about the New Right, regardless of whether or not it has anything to do with Venner's views.''

If you wanted to respectfully disagree with Venner, you ought to have shown more knowledge of his work and more knowledge of him as an intellectual. You did not indicate any familiarity with more than a single one of his works and began to make sloppy and broad-brush claims that had nothing to do with him. You were not qualified to write such a post.

No more time will be wasted on this subject.

I will continue to read your blog and appreciate the work you do.

\hfill

\texttt{John Morgan on 2013-05-28 at 09:40 said:}

I don't take issue with your interpretations of Venner or the European situation, since your opinion on your own site is your prerogative. However, once again you have made a statement that isn't accurate: ``The problem remains is that he never made clear what is our identity, except in a negative sense: our identify is not-Muslim.'' Monsieur Venner actually wrote many books in which we attempted to identify what the European identity is. One of them, ``The Clash of History,'' will soon be published in English by Arktos. One may not agree with what he wrote, but he did have an answer.

\hfill

\texttt{Izak on 2013-05-28 at 13:15 said:}

OK, I did not see Cologero's response in the previous thread.

Disregard this post.

\hfill

\texttt{August on 2013-05-28 at 16:25 said:}

I understand the perspective from which Cologero delivers his criticism, and in that light it is fair. The `geste' may be judged right away according to certain absolute standards, or in the context of what it might contribute to achieving with the passage of time. To each his own.

As for the puffing of chests in the name of Christianity that consideration of Europe's struggles engenders among some commenters here, it is tending toward the cringe-worthy. Islam, in and of itself, is no more a threat to Europe than universalist Christianity and its egalitarian tendencies -- the chief problem, and M. Venner agreed, is exceedingly alien races settling permanently and extensively across European territory, be they Muslim or otherwise, and encroaching upon our home.

Both religions do in fact have Semitic foci, Romanity nothwithstanding. Many have even argued that Christianity, being the root of the modern Western deviations, is the origin of the decadence that M. Venner reacted against. The symbolism of the location of the suicide then takes on a new, less gallant, meaning.

My duty, given intimate connections, is to remind you all of Europe's indigenous Muslims, who lived well enough and preserved their European beauty as Muslims no less, until suffering the most severe dishonor at the hands of their `Christian' nationalist brother-neighbours in recent conflicts, in a prime example of what Cologero might call the complex-ridden expression of identity in a negative sense. One can be forgiven for fearing the eventualities should kindred currents happen to motivate nascent political movements in Western Europe. Though the European spiritual minorities may be small in number, they mingle their roots with yours no less.

Finally, whatever the future holds, the following is wisdom to temper passion:

``\,'' These mighty racial tides flow from the most elemental of vital urges: self-expansion and self-preservation. Both outward thrust of expanding life and counter-thrust of threatened life are equally normal phenomena. To condemn the former as ``criminal'' and the latter as ``selfish'' is either silly or hypocritical and tends to envenom with unnecessary rancor what objective fairness might keep a candid struggle, inevitable yet alleviated by mutual comprehension and respect. This is no mere plea for ``sportsmanship''; it is a very practical matter. There are critical times ahead; times in which intense race-pressures will engender high tensions and perhaps wars. If men will keep open minds and will eschew the temptation to regard those opposing their desires to defend or possess respectively as impious fiends, the struggles will lose half their bitterness, and the wars (if wars there must be) will be shorn of half their ferocity. ``\,''

-Lothrop Stoddard in `The Rising Tide of Colour'

\hfill

\texttt{Pickman on 2013-05-28 at 18:10 said:}

You didn't see it because it's a reaction to what you've posted above which is perfectly legitimate. You should stand by your statements.

Muslims and Islam are as much a physical threat as a spiritual one and there is little reason to believe otherwise no matter how many people wish to call for restraint or rationalize the current destruction by invasion of Europe. True it's not just Islam, as there are multiple non-white groups and counter-initiatic forces of equal malignity yet let us not lose sight of any urgent threats.

When another race does harm to indigenous Europeans, it's somehow our fault to the extent we daily exalt in pathological entho-masochist tendencies that is akin to a mental-aids syndrome where we have lost the healthy immune-response to a deadly infectious attack. On the contrary it has become fashionable to welcome the festering malaise and denounce those who wish to fight it.

Democratic states and your comfortable existence are held together by violence and force -- there is always a certain unavoidable necessity of needing a military, borders, policemen and rule of law even if they do often prove ineffectual. Take these away, make people hungry and uncomfortable, and any misaligned altruisms will fade away immediately.

\hfill

\texttt{Izak on 2013-05-30 at 16:53 said:}

Oh, I'm not taking anything back! Apparently I didn't check the chronology well enough.

Anyhow I have no interest in debating Islam and its effect on Europe. I basically see the attitude of Whites as something that should probably be solved with a technique akin to addiction recovery, or something like that. First step, admit there's a problem. Second step, start to recognize a higher power. Etc. I am not terribly impressed by the ``genocide'' rhetoric going on. It is cute and makes for some amusing youtube videos about double-standards or whatever, but it doesn't address the bigger problems. Debates about immigration policy are also meaningless, because the die has already been cast. If all immigration were to be curtailed, whites still do not produce enough children to save themselves. Sometimes you can take lessons from other people and demonstrate your new knowledge competitively. I think that's why Evola argued that the first couple Crusades were actually good for both religions.


\end{sffamily}
\end{footnotesize}

